% mazzi: 05
\chapter{Le forme di partenariato pubblico privato}

\section{La nozione di partenariato pubblico privato}

\textbf{Partenariato pubblico privato} (PPP): particolari forme di partecipazione
dell'\textbf{iniziativa imprenditoriale privata} finalizzate a progettazione, costruzione,
finanziamento, gestione e manutenzione di \textbf{opere pubbliche o di pubblica utilità}, ovvero
alla fornitura di beni o all'erogazione di servizi. È l'\textbf{unione degli interessi della
collettività} con quelli dell'iniziativa imprenditoriale privata.

\rubrica{Perché se ne ricorre}

La gestione di un impianto sportivo rientra fra i \textbf{servizi pubblici}, in particolare fra i
servizi sociali. Il PPP entra in gioco quando l'ente proprietario --- oltre a non voler gestire
direttamente l'impianto --- non ha o non intende impegnare risorse pubbliche per costruire nuovi
impianti o anche solo per riqualificare, riconvertire, ammodernare o mettere in sicurezza quelli
esistenti:

\begin{itemize}
  \item pesano i \textbf{vincoli di finanza pubblica}: prima il patto di stabilità, oggi il
    pareggio di bilancio;
  \item il capitale privato (società sportive, investitori, banche, imprese di costruzioni) si
    unisce in progetti che devono comunque restare volti all'\textbf{interesse pubblico}.
\end{itemize}

\rubrica{La definizione contrattuale --- art.\ 3 D.Lgs.\ 50/2016}

``Il \textbf{contratto a titolo oneroso} stipulato per iscritto con il quale una o più
\textbf{stazioni appaltanti} conferiscono a uno o più \textbf{operatori economici} per un periodo
determinato in funzione della durata dell'ammortamento dell'investimento o delle modalità di
finanziamento fissate, un complesso di attività consistenti nella \textbf{realizzazione,
trasformazione, manutenzione e gestione operativa} di un'opera in cambio della sua
\textbf{disponibilità}, o del suo \textbf{sfruttamento economico}, o della \textbf{fornitura di un
servizio} connesso all'utilizzo dell'opera stessa, con \textbf{assunzione di rischio} secondo
modalità individuate nel contratto, da parte dell'operatore.''

Gli elementi della definizione:

\begin{itemize}
  \item \textbf{titolo oneroso}: non esistono forme di partenariato a titolo gratuito. È sempre un
    \textit{do ut des} --- si realizza l'opera pubblica avendo come contropartita la possibilità di
    gestirla e sfruttarla economicamente per il periodo di concessione;
  \item \textbf{le parti}: la stazione appaltante è la parte pubblica, nella maggior parte dei casi
    il Comune committente; l'operatore economico è il soggetto affidatario e gestore;
  \item \textbf{la durata}: commisurata all'ammortamento dell'investimento o alle modalità di
    finanziamento fissate;
  \item \textbf{l'assunzione del rischio} da parte dell'operatore, secondo le modalità del
    contratto.
\end{itemize}

Realizzazione, trasformazione, manutenzione e gestione operativa in cambio della disponibilità o
dello sfruttamento economico dell'opera si adattano bene alla gestione di un impianto sportivo e
di uno stadio di calcio.

% ---------------------------------------------------------------------------

\section{L'allocazione dei rischi}

\rubrica{La qualificazione dei contratti di PPP --- art.\ 180 D.Lgs.\ 50/2016}

Il nuovo codice, per meglio qualificare i contratti di PPP, chiarisce che i \textbf{ricavi di
gestione} devono derivare:

\begin{itemize}
  \item dal \textbf{canone} riconosciuto dall'ente concedente;
  \item e/o da \textbf{ulteriori introiti} derivanti dalla gestione del servizio;
\end{itemize}

purché permanga in capo all'operatore privato, oltre al \textbf{rischio di costruzione}, anche il
\textbf{rischio di disponibilità} o --- nei casi di attività redditizia verso l'esterno --- il
\textbf{rischio della domanda}.

\rubrica{Le quattro fattispecie di rischio --- art.\ 3 D.Lgs.\ 50/2016}

\begin{itemize}
  \item \textbf{rischio operativo}: il rischio legato alla gestione dei lavori o dei servizi sul
    lato della domanda o sul lato dell'offerta o di entrambi, trasferito all'operatore economico
    nei casi dell'art.\ 180. Si considera assunto quando, in \textbf{condizioni operative normali}
    --- cioè in assenza di eventi non prevedibili --- \textbf{non è garantito il recupero degli
    investimenti} effettuati o dei costi sostenuti per la gestione. La parte di rischio trasferita
    deve comportare una \textbf{reale esposizione alle fluttuazioni del mercato}, tale per cui ogni
    potenziale perdita stimata non sia puramente nominale o trascurabile;
  \item \textbf{rischio di costruzione}: il rischio legato al ritardo nei tempi di consegna, al
    mancato rispetto degli standard di progetto, all'aumento dei costi, a inconvenienti di tipo
    tecnico nell'opera e al mancato completamento dell'opera --- è ciò che in Italia ha bloccato o
    ritardato in modo patologico molte opere pubbliche;
  \item \textbf{rischio di disponibilità}: il rischio legato alla capacità del concessionario di
    erogare le prestazioni contrattuali pattuite, sia per volume sia per standard di qualità
    previsti;
  \item \textbf{rischio di domanda}: il rischio legato ai diversi volumi di domanda del servizio
    che il concessionario deve soddisfare --- la risposta dell'utenza --- ovvero il rischio legato
    alla mancanza di utenza e quindi di flussi di cassa sufficienti a remunerare l'investimento.
\end{itemize}

Il rischio operativo deve in ogni caso restare in capo all'operatore economico privato in
\textbf{misura maggioritaria}.

% ---------------------------------------------------------------------------

\section{L'equilibrio economico-finanziario e il prezzo}

\textbf{Equilibrio economico e finanziario} (art.\ 3 comma 1): la contemporanea presenza delle
condizioni di convenienza economica e sostenibilità finanziaria.

\begin{itemize}
  \item \textbf{convenienza economica}: capacità del progetto di creare valore nell'arco
    dell'efficacia del contratto e di generare un livello di \textbf{redditività adeguato per il
    capitale investito};
  \item \textbf{sostenibilità finanziaria}: capacità del progetto di generare \textbf{flussi di
    cassa} sufficienti a garantire il \textbf{rimborso del finanziamento}.
\end{itemize}

\rubrica{Il prezzo e il limite del 49\%}

L'equilibrio economico-finanziario è il \textbf{presupposto per la corretta allocazione dei
rischi}. A tal fine, in sede di gara, l'amministrazione aggiudicatrice può stabilire anche un
\textbf{prezzo} --- più semplicemente, un contributo --- consistente in:

\begin{itemize}
  \item un \textbf{contributo pubblico};
  \item ovvero la \textbf{cessione di beni immobili} che non assolvono più a funzioni di interesse
    pubblico (è il caso, molto discusso, delle caserme in disuso).
\end{itemize}

In ogni caso il riconoscimento del prezzo, \textbf{sommato} al valore di eventuali \textbf{garanzie
pubbliche} o di ulteriori \textbf{meccanismi di finanziamento} a carico della pubblica
amministrazione, non può superare il \textbf{49\% del costo dell'investimento complessivo},
comprensivo di eventuali oneri finanziari.

\begin{itemize}
  \item la precedente formulazione della norma fissava il limite al \textbf{30\%}: ha operato dal
    momento di entrata in vigore del codice, aprile 2016, fino ad \textbf{aprile 2017};
  \item in quell'anno gli interventi in forma di PPP sono \textbf{calati fortemente}. Un impianto
    sportivo fatica a remunerare l'investimento --- con una partita ogni due settimane, per dirla
    brutalmente --- e la sua rilevanza economica, e con essa la capacità di generare flussi di
    cassa, può essere limitata: un contributo pubblico è pressoché necessario;
  \item l'innalzamento al 49\% sfrutta l'estensione più ampia consentita dalla \textbf{direttiva
    comunitaria}, che richiede il permanere del rischio in misura maggioritaria a carico del
    privato: il limite sottintende cioè che almeno il \textbf{51\% del rischio}, anche economico,
    resti in capo al soggetto privato.
\end{itemize}

% ---------------------------------------------------------------------------

\section{La contabilizzazione off balance e il ricorso al mutuo}

\rubrica{Le decisioni Eurostat}

Agli enti locali che investono in un contratto di partenariato ex art.\ 180 sono pacificamente
applicabili le \textbf{decisioni Eurostat}, che consentono di contabilizzare \textbf{\textit{off
balance}} --- fuori bilancio --- gli \textit{asset} oggetto dell'operazione, a condizione che vi
sia una \textbf{chiara allocazione del rischio di disponibilità} in capo alla controparte privata.

\begin{itemize}
  \item le risorse impiegate non gravano quindi sul bilancio dell'ente né sui suoi equilibri;
  \item la PA può così realizzare \textbf{opere strategiche} senza sottostare ai pesanti vincoli
    imposti dai principi contabili in tema di investimenti --- vincoli che, nella suddivisione del
    bilancio pubblico fra \textbf{spese correnti} e \textbf{spese di investimento}, bloccano spesso
    l'investimento diretto della parte pubblica.
\end{itemize}

\rubrica{Il mutuo per finanziare il prezzo}

Gli enti locali possono assicurare al concessionario l'equilibrio economico-finanziario degli
investimenti e della gestione anche mediante un ``prezzo''. Per reperire le risorse necessarie ai
\textbf{trasferimenti in conto capitale} a favore dei concessionari di lavori pubblici (artt.\ 165
e 180 D.Lgs.\ 50/2016) a titolo di prezzo si può \textbf{contrarre un mutuo} --- con la Cassa
Depositi e Prestiti, con l'Istituto per il Credito Sportivo e non solo --- poiché tali
trasferimenti costituiscono investimenti per i quali agli enti locali è consentito indebitarsi.

\textbf{Art.\ 3 comma 18 L.\ 350/2003} (finanziaria per il 2004, antenata dell'odierna legge di
bilancio): ai fini dell'art.\ 119 sesto comma della Costituzione, costituiscono investimenti per i
quali gli enti locali possono ricorrere all'indebitamento, fra gli altri, alla lettera h):

\begin{quote}
i \textbf{contributi agli investimenti} e i \textbf{trasferimenti in conto capitale} a seguito di
escussione delle garanzie in favore di soggetti concessionari di lavori pubblici, o di proprietari
o gestori di impianti, di reti o di dotazioni funzionali all'erogazione di servizi pubblici, o di
soggetti che erogano servizi pubblici, le cui concessioni o contratti di servizio prevedono la
\textbf{retrocessione degli investimenti} agli enti committenti alla loro scadenza, anche
anticipata.
\end{quote}

% ---------------------------------------------------------------------------

\section{Le questioni interpretative sul limite del 49\%}

Il limite introdotto dal D.Lgs.\ 50/2016 apre una serie di dubbi, non ancora risolti, su che cosa
entri nei due termini del rapporto.

\rubrica{Che cosa entra nel costo dell'investimento}

\begin{itemize}
  \item gli \textbf{interessi del mutuo} vanno computati al \textbf{valore nominale} o
    \textbf{attualizzati}? Su un intervento di 20, 25 o 30 anni la differenza non è di poco conto.
    E se vanno attualizzati, con quale tasso --- quello usato per determinare il valore degli aiuti
    di Stato \textit{de minimis}?
  \item se per la determinazione del prezzo si considerano i contributi in conto gestione, vanno
    correlativamente inserite nel costo dell'investimento anche le \textbf{spese di gestione}? Si
    pensi alle concessioni e ai contratti di partenariato che prevedono espressamente la gestione
    dell'impianto. Quali, entro quali limiti, attualizzate a quale tasso?
\end{itemize}

\rubrica{Che cosa entra nel prezzo}

\begin{itemize}
  \item i \textbf{contributi pluriennali per la gestione} --- quelli che la PA conferisce ogni anno
    per garantire l'equilibrio economico-finanziario, per esempio quando ha imposto al gestore
    \textbf{tariffe calmierate}. Vanno computati solo se non sottoposti nell'erogazione a
    condizioni (disponibilità di bilancio del Comune, calcolo come differenza fra costi sostenuti e
    ricavi della gestione)? Le condizioni sempre previste per la \textbf{revoca della concessione}
    --- non specificamente riferite ai contributi --- bastano a escluderli dal computo? E se vanno
    considerati, vanno attualizzati e a quale tasso?
  \item i contributi esclusi dal computo vanno \textbf{reinseriti} se formano oggetto di
    \textbf{cessione all'ICS} con funzione di garanzia e di pagamento del mutuo?
  \item la \textbf{concessione di diritti reali} --- tipicamente il \textbf{diritto di superficie}
    concesso al mutuatario per realizzare l'opera e per poter concedere ipoteca alla banca --- va
    computata nel prezzo, e come? Come \textbf{differenziale} fra il costo di mercato per la
    fruizione del diritto e quello effettivamente corrisposto (area che vale 10 e ceduta a 7:
    contributo pubblico pari a 3)? Il beneficio va attualizzato e a quale tasso? La questione ha
    peso perché in passato il diritto di superficie veniva concesso a \textbf{cifre simboliche},
    anche di un euro, mentre l'esistenza di un tetto al contributo pubblico impone di valorizzare
    adeguatamente anche la cessione dell'area su cui edificare;
  \item fra le \textbf{garanzie pubbliche} va considerata la \textbf{fideiussione comunale} --- oggi
    sempre più rara, sostituita dalla concessione di diritti reali? E al valore nominale o come
    \textbf{Equivalente Sovvenzione Lordo} (ESL), eventualmente calcolato col metodo fissato per gli
    aiuti di Stato?
\end{itemize}

\rubrica{L'intreccio con gli aiuti di Stato}

Gli aiuti di Stato sono in linea di massima \textbf{vietati} a livello europeo e costantemente
attenzionati dalla Commissione. Il problema dei contributi pluriennali del Comune al concessionario
rileva anche su questo piano: si devono considerare solo quando sono \textbf{trasparenti}, cioè
calcolabili con precisione \textit{ex ante}, al momento della stipula. Ci si chiede allora se
occorra tener conto della possibile revoca della concessione per inadempimenti contrattuali ---
escludendoli dal computo --- oppure se vadano esclusi solo in presenza di clausole specificamente
riferite ai contributi che ne rendano incerta l'erogazione.

\rubrica{La ricostruzione proposta}

Prevale una lettura \textbf{non restrittiva}, coerente con la direttiva europea:

\begin{itemize}
  \item nel \textbf{costo dell'investimento} su cui calcolare il 49\% vanno inseriti il costo
    totale delle opere, il costo del finanziamento comprensivo di \textbf{interessi attualizzati} e
    gli \textbf{oneri di gestione complessivi};
  \item nel \textbf{prezzo}, unito agli altri ``aiuti pubblici'', vanno sommati:
    \begin{itemize}
      \item trasferimenti in conto capitale e immobili;
      \item valore dei \textbf{diritti reali di godimento} (p.\,es.\ il diritto di superficie);
      \item \textbf{contributi pluriennali} anche in conto gestione, attualizzati e anche se
        sottoposti a condizioni;
      \item valore della \textbf{fideiussione comunale} e di altre garanzie pubbliche, calcolato
        come ESL;
      \item valore attualizzato di altri \textbf{contributi pubblici pluriennali}, compresi quelli
        in conto interessi.
    \end{itemize}
\end{itemize}

% ---------------------------------------------------------------------------

\section{I fondi immobiliari}

\textbf{Art.\ 33 comma 2 D.L.\ 6 luglio 2011, n.\ 98}: i comuni privi di risorse proprie da
destinare al recupero, alla ristrutturazione e alla valorizzazione di impianti sportivi --- stadi,
palazzetti dello sport --- possono \textbf{apportare} detti impianti, o i diritti reali a essi
relativi, a un \textbf{fondo immobiliare} costituito \textit{ad hoc}.

\begin{itemize}
  \item in cambio gli enti ricevono \textbf{quote del fondo};
  \item il fondo trae dalla \textbf{valorizzazione degli \textit{asset} immobiliari} conferiti la
    remunerazione dei capitali raccolti sul mercato e in esso investiti.
\end{itemize}

% ---------------------------------------------------------------------------

\section{Le tipologie di contratti di PPP}

Nella tipologia dei contratti di PPP rientrano:

\begin{itemize}
  \item la \textbf{finanza di progetto} (\textit{project financing}), in cui la fase progettuale è
    a carico del privato;
  \item la \textbf{concessione di costruzione e gestione}, che dal \textit{project financing}
    differisce perché è la parte pubblica ad accollarsi gli \textbf{oneri di progettazione}: il
    capitolato e le caratteristiche strutturali dell'impianto vanno a bando con studi progettuali
    già approvati dall'amministrazione;
  \item la \textbf{concessione di servizi}, dove non vi è necessariamente una componente
    preponderante di lavori e la remunerazione viene dalla prevalente gestione del servizio
    pubblico;
  \item la \textbf{locazione finanziaria di opere pubbliche} (\textit{leasing});
  \item il \textbf{contratto di disponibilità};
\end{itemize}

e \textbf{qualunque altra procedura} di realizzazione in partenariato di opere o servizi che
presenti le caratteristiche dei commi precedenti --- formula sufficientemente ampia da comprendere
anche schemi non ancora codificati.

% ---------------------------------------------------------------------------

\section{Le procedure di affidamento}

La scelta dell'operatore economico avviene con \textbf{procedure ad evidenza pubblica}, anche
mediante \textbf{dialogo competitivo}.

Salva l'ipotesi in cui l'affidamento abbia a oggetto \textbf{anche l'attività di progettazione}, le
amministrazioni aggiudicatrici pongono a base di gara:

\begin{itemize}
  \item il \textbf{progetto definitivo};
  \item uno \textbf{schema di contratto} --- la convenzione per la gestione dell'impianto e il
    regolamento d'uso;
  \item un \textbf{piano economico finanziario};
\end{itemize}

che disciplinino l'\textbf{allocazione dei rischi} fra amministrazione aggiudicatrice e operatore
economico, verificando che il rischio permanga in misura maggioritaria --- almeno il 51\% --- a
carico dell'operatore privato.
